\scanpage{22}
\uline{Pb}\qquad $X$ sur $S$ en car $p>0$

a) Définir une opération ``tensorielle''
$\mathcal O_X$-linéaire
\[
\underline\Omega^1_{X/S}\longrightarrow
?\deleted{\unsure{\operatorname{Symm}}^p}(\Omega^1_{X/S})
\]
qui soit transposée dans un sens convenable
de \deleted{$\unsure{\partial}$} $\partial\mathbin{\text{---}}\partial^p$
sur les $S$-dérivations, et compatible avec les carrés
\diagram{\begin{tikzcd}
X \arrow[d] & X' \arrow[l] \arrow[d]\\
S & S' \arrow[l]
\end{tikzcd}}

Il faudrait que lorsque $\omega$ est une forme \uline{fermée},
on trouve quelque chose dans
$\deleted{\unsure{\operatorname{Symm}}^p_{\underline O_X}}\boxed{\underline\Omega^1_{X/S}{}^{(p)}}$
(puissance Frobéniusienne de $\underline\Omega^1_{X/S}$),
ce qui est l'opération de Cartier.
Si les $\omega_i$, $\partial_j$ sont tels que
\unsure{on n'en sort pas} par $d$, $[\ ,\ ]$, et que
les produits scalaires $\langle\omega_i,\partial_j\rangle$ sont des c\textsuperscript{tes},
on voit qu'on doit \unsure{échouer|aboutir} dans
$\boxed{\unsure{\operatorname{Symm}^p(\underline\Omega^1_{X/S})|\operatorname{Symm}^p(\underline\Omega^1_{X/S})^p}}$.

b) Montrer qu'un groupe infinitésimal d'ordre $1$ sur $S$
est défini par $\Omega$ qu. coh. sur $S$, avec
\[
C:\Omega\longrightarrow\operatorname{Symm}_{\underline O_S}^p(\Omega)
\]
\[
d_1:\Omega\mathbin{\text{---}}\bigwedge_{\underline O_S}^2\Omega
\]
\margin{\unsure{satisfaisant} certaines équations.}
\margin{Dire que $G$ est commutatif signifie que $d_1=0$.}
\margin{c) Définir le ``Verschiebung'' $V_S(X):p_S(G)\mathbin{\text{---}}G$
pour un schéma en groupes comm. $G$ sur $S$.}
\mspage{21}
